%*************************************
% ภาคผนวก ง. (ต่อ) - กระบวนการทำงานของ Desktop Agent
%*************************************

\section{ภาพรวม Desktop Agent}
\hspace*{1.3em}
Desktop Agent เป็นโปรแกรมที่ทำงานบนเครื่องคอมพิวเตอร์ของผู้ใช้ ทำหน้าที่เป็นตัวกลาง (Relay) ระหว่าง Backend บน Vercel Serverless กับอุปกรณ์เครือข่ายภายในองค์กร เนื่องจาก Vercel Serverless Functions ไม่สามารถเชื่อมต่อ SSH/NETCONF ไปยังอุปกรณ์ในเครือข่ายภายในได้โดยตรง Agent จึงทำหน้าที่รับคำสั่งจาก Backend ผ่าน HTTP Polling ดำเนินการกับอุปกรณ์เครือข่ายในเครื่อง และส่งผลลัพธ์กลับไปยัง Backend

\hspace*{1.3em}
Agent ประกอบด้วย 4 Handler หลัก ดังนี้

\begin{mycustomenum2}
    \item \textbf{SSH Handler:} จัดการการเชื่อมต่อ SSH ไปยังอุปกรณ์เครือข่าย สำหรับสั่งงาน CLI, สำรองข้อมูล (Backup) และ Deploy Configuration
    \item \textbf{NETCONF Handler:} จัดการ NETCONF Session ผ่าน SSH Subsystem สำหรับ Deploy Configuration แบบ XML/YANG
    \item \textbf{Console Handler:} จัดการการเชื่อมต่อผ่าน Serial Port (Console Cable) สำหรับตั้งค่าเริ่มต้นอุปกรณ์ที่ยังไม่มี IP Address
    \item \textbf{Ollama Handler:} เชื่อมต่อกับ Ollama LLM Server ที่ทำงานบนเครื่องผู้ใช้ สำหรับสร้าง Configuration จากภาษาธรรมชาติ
\end{mycustomenum2}

\section{การสื่อสารแบบ HTTP Polling}
\hspace*{1.3em}
เนื่องจาก Backend อยู่บน Vercel Serverless ซึ่งไม่รองรับ WebSocket (Persistent Connection) ระบบจึงใช้ HTTP Polling เป็นกลไกการสื่อสารระหว่าง Agent กับ Backend โดย Agent จะส่ง HTTP Request ไปยัง Backend เป็นระยะ ๆ เพื่อ ตรวจสอบคำสั่งใหม่ และรายงานสถานะ

\hspace*{1.3em}
ตารางที่ \ref{tabD2:PollingEndpoints} แสดง API Endpoints ที่ Agent ใช้ในการสื่อสารกับ Backend

\begin{table}[H]
\centering
\caption{\fontSixTeen{API Endpoints สำหรับ Agent Polling}}
\label{tabD2:PollingEndpoints}
\fontSixTeen
\begin{tabularx}{\textwidth}{|l|l|l|X|}
\hline
\textbf{Endpoint} & \textbf{Method} & \textbf{Interval} & \textbf{คำอธิบาย} \\
\hline
/api/agent/poll/heartbeat & POST & 5 วินาที & ส่ง Heartbeat แจ้งสถานะ Online พร้อมข้อมูล Serial Port \\
\hline
/api/agent/poll/commands & GET & 2 วินาที & ดึงคำสั่งที่รอดำเนินการจาก Command Queue \\
\hline
/api/agent/poll/result/\{id\} & POST & ทุกคำสั่ง & ส่งผลลัพธ์การดำเนินการกลับไปยัง Backend \\
\hline
/api/agent/poll/shell-create & POST & ครั้งเดียว & สร้าง Interactive Shell Session ใหม่ \\
\hline
/api/agent/poll/shell-output & POST & ทุก Output & ส่ง Output จาก Shell ไปยัง Frontend \\
\hline
/api/agent/poll/shell-input/\{id\} & GET & 500 ms & ดึง Keystroke Input จาก Frontend \\
\hline
/api/agent/poll/shell-closed & POST & ครั้งเดียว & แจ้ง Backend ว่า Shell Session ถูกปิด \\
\hline
\end{tabularx}
\end{table}

\hspace*{1.3em}
ทุก HTTP Request ที่ Agent ส่งจะแนบ Header \textbf{X-Agent-Token} สำหรับยืนยันตัวตน โดย Token นี้ได้มาจากการลงทะเบียน Agent ผ่านหน้าเว็บแอปพลิเคชัน

\section{กระบวนการเริ่มต้นและเชื่อมต่อ Agent}
\hspace*{1.3em}
เมื่อ Agent เริ่มทำงาน จะมีลำดับขั้นตอนดังนี้

\begin{mycustomenum2}
    \item \textbf{โหลดการตั้งค่า:} อ่านค่า Server URL, Agent Token และ Agent Name จาก Configuration Store ที่เก็บไว้ในเครื่องผู้ใช้
    \item \textbf{ตั้งค่าเริ่มต้น (ครั้งแรก):} หากยังไม่มีการตั้งค่า ระบบจะแสดง Interactive CLI ให้ผู้ใช้กรอก Server URL และ Agent Token
    \item \textbf{ตรวจสอบ Ollama:} ตรวจสอบว่า Ollama LLM Server พร้อมใช้งานหรือไม่ หากไม่พร้อม Agent จะยังทำงานได้แต่ไม่สามารถสร้าง Configuration ด้วย AI ได้
    \item \textbf{สร้าง Handlers:} สร้าง SSH Handler, NETCONF Handler, Console Handler และ Ollama Handler
    \item \textbf{เชื่อมต่อ Backend:} ส่ง Heartbeat แรกเพื่อยืนยันตัวตน เมื่อสำเร็จจะเริ่ม Polling Loop สำหรับ Heartbeat (ทุก 5 วินาที) และ Command Queue (ทุก 2 วินาที)
    \item \textbf{รอรับคำสั่ง:} Agent เข้าสู่สถานะ Online พร้อมรับคำสั่งจาก Backend
\end{mycustomenum2}

\section{กระบวนการ Command Queue}
\hspace*{1.3em}
Command Queue เป็นกลไกหลักในการสื่อสารระหว่าง Backend กับ Agent โดยมีขั้นตอนดังนี้

\begin{mycustomenum2}
    \item \textbf{ผู้ใช้สั่งงาน:} ผู้ใช้ดำเนินการผ่านหน้าเว็บ เช่น สั่ง Backup, Deploy หรือ Test Connection
    \item \textbf{Backend สร้าง AgentCommand:} Backend สร้าง Document ใน MongoDB Collection \textbf{agentcommands} พร้อมสถานะ \textbf{pending} ระบุ Event, Data และ User ID
    \item \textbf{Agent ดึงคำสั่ง:} Agent ทำ HTTP GET ไปยัง \textbf{/api/agent/poll/commands} ทุก 2 วินาที Backend จะ query คำสั่งที่มีสถานะ pending และเปลี่ยนเป็น \textbf{processing}
    \item \textbf{Agent ดำเนินการ:} Agent รับคำสั่ง เรียก Handler ที่เกี่ยวข้อง (SSH, NETCONF, Console หรือ Ollama) เพื่อดำเนินการกับอุปกรณ์จริง
    \item \textbf{Agent ส่งผลลัพธ์:} Agent ทำ HTTP POST ผลลัพธ์ไปยัง \textbf{/api/agent/poll/result/\{commandId\}} Backend อัปเดตสถานะ Command เป็น \textbf{completed} หรือ \textbf{failed}
    \item \textbf{Frontend รับผล:} Frontend ทำ Polling ไปยัง Backend เพื่อดึงผลลัพธ์และแสดงต่อผู้ใช้
\end{mycustomenum2}

\section{ประเภทคำสั่งที่ Agent รองรับ}
\hspace*{1.3em}
Agent รองรับคำสั่งทั้งหมด 23 ประเภท แบ่งตาม Handler ดังตารางที่ \ref{tabD2:AgentCommands}

\begin{table}[H]
\centering
\caption{\fontSixTeen{คำสั่ง SSH ที่ Agent รองรับ}}
\label{tabD2:AgentCommands}
\fontSixTeen
\begin{tabularx}{\textwidth}{|l|X|}
\hline
\textbf{Event} & \textbf{คำอธิบาย} \\
\hline
agent:ssh:connect & เชื่อมต่อ SSH ไปยังอุปกรณ์เครือข่าย \\
\hline
agent:ssh:disconnect & ปิดการเชื่อมต่อ SSH \\
\hline
agent:ssh:exec & สั่งคำสั่ง CLI เดี่ยวบนอุปกรณ์ \\
\hline
agent:ssh:send-config & ส่งชุดคำสั่ง Configuration ผ่าน Interactive Shell \\
\hline
agent:ssh:deploy-config & เชื่อมต่อ $\rightarrow$ ส่ง Config $\rightarrow$ ปิดการเชื่อมต่อ (All-in-one) \\
\hline
agent:ssh:backup & เชื่อมต่อ $\rightarrow$ ดึง Running/Startup Config $\rightarrow$ ปิด (All-in-one) \\
\hline
agent:ssh:open-shell & เปิด Interactive SSH Shell Session \\
\hline
agent:ssh:shell-input & ส่ง Keystroke ไปยัง Shell Session ที่เปิดอยู่ \\
\hline
agent:ssh:close-shell & ปิด Interactive Shell Session \\
\hline
\end{tabularx}
\end{table}

\begin{table}[H]
\centering
\caption{\fontSixTeen{คำสั่ง NETCONF ที่ Agent รองรับ}}
\label{tabD2:NetconfCommands}
\fontSixTeen
\begin{tabularx}{\textwidth}{|l|X|}
\hline
\textbf{Event} & \textbf{คำอธิบาย} \\
\hline
agent:netconf:connect & สร้าง NETCONF Session (SSH Subsystem Port 830) \\
\hline
agent:netconf:disconnect & ปิด NETCONF Session \\
\hline
agent:netconf:get-config & อ่าน Running Configuration ผ่าน get-config RPC \\
\hline
agent:netconf:edit-config & แก้ไข Configuration ผ่าน Candidate Datastore Workflow \\
\hline
agent:netconf:rpc & ส่ง Raw RPC Operation \\
\hline
\end{tabularx}
\end{table}

\begin{table}[H]
\centering
\caption{\fontSixTeen{คำสั่ง Console (Serial Port) ที่ Agent รองรับ}}
\label{tabD2:ConsoleCommands}
\fontSixTeen
\begin{tabularx}{\textwidth}{|l|X|}
\hline
\textbf{Event} & \textbf{คำอธิบาย} \\
\hline
agent:console:list-ports & แสดงรายการ Serial Port ที่เชื่อมต่อกับเครื่อง \\
\hline
agent:console:connect & เชื่อมต่อ Serial Port \\
\hline
agent:console:disconnect & ปิดการเชื่อมต่อ Serial Port \\
\hline
agent:console:command & ส่งคำสั่งผ่าน Serial Console \\
\hline
agent:console:test & ทดสอบการเชื่อมต่อ Serial Port \\
\hline
agent:console:initial-config & ส่งชุดคำสั่งตั้งค่าเริ่มต้น SSH ให้อุปกรณ์ \\
\hline
\end{tabularx}
\end{table}

\begin{table}[H]
\centering
\caption{\fontSixTeen{คำสั่ง Ollama LLM ที่ Agent รองรับ}}
\label{tabD2:OllamaCommands}
\fontSixTeen
\begin{tabularx}{\textwidth}{|l|X|}
\hline
\textbf{Event} & \textbf{คำอธิบาย} \\
\hline
agent:ollama:status & ตรวจสอบสถานะ Ollama Server \\
\hline
agent:ollama:health & ตรวจสอบว่า Ollama พร้อมใช้งานหรือไม่ \\
\hline
agent:ollama:models & แสดงรายการ LLM Model ที่ติดตั้ง \\
\hline
agent:ollama:set-model & เปลี่ยน Model ที่ใช้งาน \\
\hline
agent:ollama:chat & สร้าง Configuration ด้วย Chat Completion \\
\hline
agent:ollama:generate & สร้างข้อความด้วย Text Generation \\
\hline
\end{tabularx}
\end{table}

\section{กระบวนการ Heartbeat}
\hspace*{1.3em}
Agent ส่ง Heartbeat ไปยัง Backend ทุก 5 วินาทีเพื่อรายงานสถานะ Online/Offline โดย Heartbeat Payload ประกอบด้วยข้อมูลดังนี้

\begin{mycustomenum2}
    \item \textbf{agentName:} ชื่อ Agent ที่ผู้ใช้กำหนด
    \item \textbf{agentVersion:} เวอร์ชันของ Agent
    \item \textbf{platform:} ระบบปฏิบัติการ (win32, linux, darwin)
    \item \textbf{hostname:} ชื่อเครื่องคอมพิวเตอร์
    \item \textbf{capabilities.serialPorts:} รายการ Serial Port ที่เชื่อมต่อกับเครื่อง (สแกนทุกครั้งที่ส่ง Heartbeat)
\end{mycustomenum2}

\hspace*{1.3em}
Backend จะบันทึก Heartbeat ลงใน Collection \textbf{agentheartbeats} และใช้เวลาของ Heartbeat ล่าสุดในการตัดสินว่า Agent ยังออนไลน์อยู่หรือไม่ หากไม่ได้รับ Heartbeat เกิน 30 วินาที Backend จะถือว่า Agent ออฟไลน์

\section{กระบวนการ SSH Backup}
\hspace*{1.3em}
การสำรองข้อมูล Configuration ของอุปกรณ์ผ่าน Agent มีขั้นตอนดังนี้

\begin{mycustomenum2}
    \item \textbf{เชื่อมต่อ SSH:} Agent สร้าง SSH Connection ไปยังอุปกรณ์ที่ระบุ
    \item \textbf{เปิด Interactive Shell:} เปิด Shell Session เพื่อรองรับ Multi-command Sequence บนอุปกรณ์ Cisco
    \item \textbf{ปิด Pagination:} ส่งคำสั่ง \textbf{terminal length 0} เพื่อปิดการแบ่งหน้า Output
    \item \textbf{ดึง Configuration:} ส่งคำสั่ง \textbf{show running-config} และ/หรือ \textbf{show startup-config} ตามประเภทที่เลือก โดยใช้ Shell Session เดียวกันเพื่อหลีกเลี่ยงปัญหา VTY Line
    \item \textbf{ทำความสะอาด Output:} ลบ ANSI Escape Codes, Command Echo และ Device Prompt ออกจาก Output เพื่อให้ได้เฉพาะ Configuration จริง
    \item \textbf{ปิดการเชื่อมต่อ:} ปิด SSH Connection
    \item \textbf{ส่งผลลัพธ์:} ส่ง Running Config และ/หรือ Startup Config กลับไปยัง Backend เพื่อบันทึกลงฐานข้อมูล
\end{mycustomenum2}

\hspace*{1.3em}
ระบบรองรับ 3 โหมดการ Backup ดังตารางที่ \ref{tabD2:BackupModes}

\begin{table}[H]
\centering
\caption{\fontSixTeen{โหมดการ Backup Configuration}}
\label{tabD2:BackupModes}
\fontSixTeen
\begin{tabularx}{\textwidth}{|l|l|X|}
\hline
\textbf{โหมด} & \textbf{คำสั่ง SSH} & \textbf{คำอธิบาย} \\
\hline
running-config & show running-config & ดึงเฉพาะ Running Configuration \\
\hline
startup-config & show startup-config & ดึงเฉพาะ Startup Configuration \\
\hline
both & show running-config + show startup-config & ดึงทั้งสองใน Shell Session เดียว \\
\hline
\end{tabularx}
\end{table}

\section{กระบวนการ Deploy Configuration ผ่าน SSH}
\hspace*{1.3em}
การ Deploy Configuration ผ่าน SSH จะใช้ Interactive Shell เพื่อส่งชุดคำสั่ง CLI ไปยังอุปกรณ์ โดยมีลำดับขั้นตอนดังนี้

\begin{mycustomenum2}
    \item \textbf{เชื่อมต่อ SSH:} Agent สร้าง SSH Connection ไปยังอุปกรณ์เป้าหมาย
    \item \textbf{เปิด Shell:} เปิด Interactive Shell พร้อม Terminal Emulation (xterm)
    \item \textbf{เข้าสู่ Privileged Mode:} ส่งคำสั่ง \textbf{enable} พร้อม Enable Password (ถ้ามี)
    \item \textbf{เข้าสู่ Configuration Mode:} ส่งคำสั่ง \textbf{configure terminal}
    \item \textbf{ส่งคำสั่ง Config:} ส่งคำสั่ง Configuration ทีละบรรทัด
    \item \textbf{ออกจาก Configuration Mode:} ส่งคำสั่ง \textbf{end}
    \item \textbf{บันทึก Configuration:} ส่งคำสั่ง \textbf{write memory} เพื่อบันทึกลง Startup Config
\end{mycustomenum2}

\hspace*{1.3em}
ระบบใช้ Prompt Detection Pattern \textbf{/[\#>]\textbackslash s*\$/} พร้อม Debounce 200ms เพื่อตรวจจับว่าอุปกรณ์พร้อมรับคำสั่งถัดไปแล้ว โดยมี Timeout สูงสุด 60 วินาที

\section{กระบวนการ Interactive Shell ผ่าน Agent Relay}
\hspace*{1.3em}
Interactive Shell เป็นฟีเจอร์ที่ช่วยให้ผู้ใช้สามารถใช้งาน SSH Terminal ผ่านเว็บเบราว์เซอร์ได้ โดย Agent ทำหน้าที่ Relay I/O ระหว่าง Frontend กับอุปกรณ์เครือข่ายผ่าน Polling-based Protocol ดังนี้

\begin{mycustomenum2}
    \item \textbf{สร้าง Session:} Agent สร้าง Shell Session บนอุปกรณ์และแจ้ง Backend ผ่าน \textbf{POST /api/agent/poll/shell-create}
    \item \textbf{ส่ง Output:} เมื่อมี Data จากอุปกรณ์ Agent จะส่ง Output ไปยัง Backend ผ่าน \textbf{POST /api/agent/poll/shell-output} ทันที
    \item \textbf{รับ Input:} Agent ทำ Polling ไปยัง \textbf{GET /api/agent/poll/shell-input/\{sessionId\}} ทุก 500ms เพื่อรับ Keystroke จากผู้ใช้
    \item \textbf{ส่ง Input ไปยังอุปกรณ์:} เมื่อได้รับ Keystroke Agent จะเขียนลง SSH Shell Stream
    \item \textbf{ปิด Session:} เมื่อ Shell ถูกปิด Agent จะแจ้ง Backend ผ่าน \textbf{POST /api/agent/poll/shell-closed}
\end{mycustomenum2}

\hspace*{1.3em}
Session ID มีรูปแบบ \textbf{shell-\{timestamp\}-\{random6\}} และข้อมูล I/O ถูกเก็บใน MongoDB Collection \textbf{shellsessions} เพื่อ Buffer ระหว่าง Agent กับ Frontend

\section{กระบวนการเชื่อมต่อ Serial Console}
\hspace*{1.3em}
Console Handler ใช้สำหรับเชื่อมต่อกับอุปกรณ์เครือข่ายผ่านสาย Console (Serial Port) เพื่อตั้งค่าเริ่มต้น เช่น กำหนด IP Address และเปิดใช้งาน SSH ให้กับอุปกรณ์ใหม่ที่ยังไม่มีการกำหนดค่า

\hspace*{1.3em}
ตารางที่ \ref{tabD2:SerialConfig} แสดงค่าเริ่มต้นของ Serial Port Configuration

\begin{table}[H]
\centering
\caption{\fontSixTeen{ค่าเริ่มต้น Serial Port Configuration}}
\label{tabD2:SerialConfig}
\fontSixTeen
\begin{tabularx}{\textwidth}{|l|l|X|}
\hline
\textbf{พารามิเตอร์} & \textbf{ค่าเริ่มต้น} & \textbf{คำอธิบาย} \\
\hline
Baud Rate & 9600 & ความเร็วในการสื่อสาร (bps) \\
\hline
Data Bits & 8 & จำนวนบิตข้อมูลต่อ Frame \\
\hline
Parity & none & การตรวจสอบความถูกต้อง \\
\hline
Stop Bits & 1 & จำนวนบิตหยุด \\
\hline
\end{tabularx}
\end{table}

\hspace*{1.3em}
ระบบจะสแกน Serial Port ทุกครั้งที่ส่ง Heartbeat และระบุ Manufacturer (FTDI, Prolific, Silicon Labs) เพื่อช่วยผู้ใช้เลือก Port ที่ถูกต้อง นอกจากนี้ Backend จะตรวจจับ \textbf{serialPortAvailable} Flag เพื่อ Relay คำสั่ง Console ไปยัง Desktop Agent โดยอัตโนมัติ เนื่องจาก Vercel Serverless ไม่สามารถใช้ Native Module ของ Serial Port ได้

\section{กระบวนการเชื่อมต่อ Ollama LLM}
\hspace*{1.3em}
Ollama Handler จัดการการเชื่อมต่อกับ Ollama LLM Server ที่ทำงานบนเครื่องผู้ใช้ (localhost:11434) สำหรับสร้าง Network Configuration จากภาษาธรรมชาติ

\hspace*{1.3em}
ตารางที่ \ref{tabD2:OllamaConfig} แสดงค่าเริ่มต้นของ Ollama Configuration

\begin{table}[H]
\centering
\caption{\fontSixTeen{ค่าเริ่มต้น Ollama LLM Configuration}}
\label{tabD2:OllamaConfig}
\fontSixTeen
\begin{tabularx}{\textwidth}{|l|l|X|}
\hline
\textbf{พารามิเตอร์} & \textbf{ค่าเริ่มต้น} & \textbf{คำอธิบาย} \\
\hline
Base URL & http://localhost:11434 & URL ของ Ollama Server \\
\hline
Default Model & qwen2.5-coder & โมเดล LLM ที่ใช้เป็นค่าเริ่มต้น \\
\hline
Request Timeout & 120 วินาที & Timeout สำหรับ API Request \\
\hline
Temperature & 0.1 & ค่าความสุ่มของ Output (ต่ำ = แม่นยำ) \\
\hline
Max Tokens & 1000 & จำนวน Token สูงสุดที่สร้าง \\
\hline
Top-p & 0.85 & Nucleus Sampling Parameter \\
\hline
\end{tabularx}
\end{table}

\hspace*{1.3em}
Ollama Handler จะแปลง Response จาก Ollama API ให้อยู่ในรูปแบบ OpenAI-compatible เพื่อให้ Backend สามารถใช้งานได้โดยไม่ต้องแก้ไข Logic โดยรูปแบบ Response ที่ส่งกลับมีโครงสร้างดังนี้

\begin{verbatim}
{
  "choices": [{
    "message": { "role": "assistant", "content": "..." },
    "finish_reason": "stop"
  }],
  "usage": {
    "prompt_tokens": 150,
    "completion_tokens": 200,
    "total_tokens": 350
  },
  "model": "qwen2.5-coder",
  "provider": "ollama"
}
\end{verbatim}

\section{ค่าเริ่มต้นและ Timeout ของ Agent}
\hspace*{1.3em}
ตารางที่ \ref{tabD2:AgentTimeouts} สรุปค่า Timeout และค่าเริ่มต้นที่สำคัญของแต่ละ Handler

\begin{table}[H]
\centering
\caption{\fontSixTeen{สรุปค่า Timeout และค่าเริ่มต้นของ Agent Handlers}}
\label{tabD2:AgentTimeouts}
\fontSixTeen
\begin{tabularx}{\textwidth}{|l|l|l|X|}
\hline
\textbf{Handler} & \textbf{Operation} & \textbf{Timeout} & \textbf{คำอธิบาย} \\
\hline
SSH & Connect & 10 วินาที & เชื่อมต่อ SSH \\
\hline
SSH & Execute Command & 30 วินาที & สั่งคำสั่งเดี่ยว \\
\hline
SSH & Send Config & 60 วินาที & ส่งชุดคำสั่ง Config \\
\hline
SSH & Backup (single) & 20 วินาที & ดึง Config ประเภทเดียว \\
\hline
SSH & Backup (both) & 30 วินาที & ดึงทั้ง Running และ Startup \\
\hline
SSH & Session Timeout & 10 นาที & หมดเวลา Shell Session \\
\hline
NETCONF & Connect & 30 วินาที & สร้าง SSH + NETCONF Session \\
\hline
NETCONF & Hello Fallback & 5 วินาที & Fallback ถ้าไม่ได้รับ Server Hello \\
\hline
NETCONF & Send RPC & 30 วินาที & ส่ง XML RPC Operation \\
\hline
NETCONF & Session Timeout & 5 นาที & หมดเวลา NETCONF Session \\
\hline
Console & Command (prompt) & 10 วินาที & รอ Prompt หลังส่งคำสั่ง \\
\hline
Console & Command (no wait) & 2 วินาที & ไม่รอ Prompt \\
\hline
Ollama & Request & 120 วินาที & LLM Chat Completion \\
\hline
Ollama & Health Check & 5 วินาที & ตรวจสอบ Ollama \\
\hline
Polling & Heartbeat & 5 วินาที & รอบ Heartbeat \\
\hline
Polling & Commands & 2 วินาที & รอบดึงคำสั่ง \\
\hline
Polling & Shell Input & 500 ms & รอบดึง Keystroke \\
\hline
\end{tabularx}
\end{table}

\section{SSH Algorithm ที่รองรับ}
\hspace*{1.3em}
Agent รองรับ SSH Algorithm หลากหลายเพื่อให้เข้ากันได้กับอุปกรณ์เครือข่ายทั้งรุ่นใหม่และรุ่นเก่า ดังตารางที่ \ref{tabD2:SSHAlgorithms}

\begin{table}[H]
\centering
\caption{\fontSixTeen{SSH Algorithms ที่ Agent รองรับ}}
\label{tabD2:SSHAlgorithms}
\fontSixTeen
\begin{tabularx}{\textwidth}{|l|X|}
\hline
\textbf{ประเภท} & \textbf{Algorithms ที่รองรับ} \\
\hline
Key Exchange & ecdh-sha2-nistp256, ecdh-sha2-nistp384, ecdh-sha2-nistp521, diffie-hellman-group14-sha256, diffie-hellman-group14-sha1, diffie-hellman-group1-sha1 \\
\hline
Cipher & aes128-ctr, aes192-ctr, aes256-ctr, aes128-cbc, aes192-cbc, aes256-cbc, 3des-cbc, aes128-gcm, aes256-gcm \\
\hline
HMAC & hmac-sha2-256, hmac-sha1, hmac-sha2-512 \\
\hline
Host Key & ssh-rsa, ssh-dss, ecdsa-sha2-nistp256, ecdsa-sha2-nistp384, ecdsa-sha2-nistp521, ssh-ed25519, rsa-sha2-256, rsa-sha2-512 \\
\hline
\end{tabularx}
\end{table}

\section{Dependencies ของ Agent}
\hspace*{1.3em}
ตารางที่ \ref{tabD2:AgentDeps} แสดง Package Dependencies หลักของ Agent

\begin{table}[H]
\centering
\caption{\fontSixTeen{Package Dependencies ของ Desktop Agent}}
\label{tabD2:AgentDeps}
\fontSixTeen
\begin{tabularx}{\textwidth}{|l|l|X|}
\hline
\textbf{Package} & \textbf{Version} & \textbf{หน้าที่} \\
\hline
ssh2 & \^{}1.15.0 & เชื่อมต่อ SSH และ NETCONF Session \\
\hline
serialport & \^{}12.0.0 & เชื่อมต่อ Serial Console Port \\
\hline
@serialport/parser-readline & \^{}12.0.0 & Parse ข้อมูล Serial แบบ Line-based \\
\hline
conf & \^{}12.0.0 & จัดเก็บ Configuration แบบ Persistent \\
\hline
chalk & \^{}5.3.0 & แสดงสีใน Terminal \\
\hline
ora & \^{}7.0.1 & แสดง Spinner ใน Terminal \\
\hline
inquirer & \^{}9.2.12 & Interactive CLI Prompts \\
\hline
\end{tabularx}
\end{table}

\hspace*{1.3em}
Agent ถูก Bundle ด้วย \textbf{esbuild} และ Build เป็น Standalone Executable (\textbf{NetConfigAgent.exe}) ด้วย \textbf{@yao-pkg/pkg} โดย Target เป็น \textbf{node18-win-x64} เพื่อให้ผู้ใช้สามารถรันได้โดยไม่ต้องติดตั้ง Node.js


